\vspace{-0.2cm}
\section{Introduction}
\vspace{-0.2cm}


Decoding visual stimuli from brain activity provides a unique lens into human perception \citep{naselaris2011encoding, haufe2014interpretation}. A central approach uses fMRI signals—which measure neural activity through blood-oxygen-level-dependent responses—to reconstruct the images perceived by subjects \citep{allen2022massive, chang2019bold5000, luo2023brainscuba}. Recent advances in deep generative models have significantly improved these reconstructions, deepening our understanding of visual representation in the brain \citep{chen2023seeing} and enabling applications in brain-computer interfaces \citep{sitaram2008fmri} and brain-driven content generation \citep{wang2024braindreamer, qiumindllm}.
%\yj{The first and last sentences are good. The middle ones do not connect well. You do not need to mention past work in this paragraph. It would be better that you start with the meaning of decoding human brain, e.g. why do researchers want to decode human brain, what are potential impacts? Among different tools to decode human brains, why using fmri data to generate image is a good way? Why fmri? why image generation?}

% make high level, universal domain, be abstract, concepts, we don't know what the space looks like. we then check if it closes to ML space.
% understanding way is align with cv model processing, if not? close to which kind of.use descriptive language to define. representation space align 



%\yj{The following challenges sound very specific and ill-motivated. It's not clear why we care about these challenges. My suggestions: First talk about general drawbacks of the prior works. In general, say that the prior works focus on the improving generation quality while neglecting the neuroscientific insights. list some papers. Then talk about the three challenges with proper transitions.}
% The success of image-fMRI reconstruction depends on two key factors: the quality of the transformation from fMRI signals to the generative model’s latent space (alignment) and the model’s capacity to generate high-quality images. 
FMRI-to-Image reconstruction involves two stages: mapping fMRI signals into a latent space and then generating images from that space. Its success depends on the similarity between the latent space and neural activity (alignment) and how well the generative model produces high-quality images.
%\yy{please do not use the word synthesize. use generate instead. this is the terminology used in other literature, then follow the literature. same for the remaining usages.}
While recent studies \citep{scotti2023reconstructing, scotti2024mindeye2, mai2024brain} focus on enhancing image quality using advanced generative models \citep{podell2023sdxl, xu2023versatile}, alignment remains underexplored. Prior work often assumes that the latent space should match the modality of the stimuli, i.e., using vision model representations to reconstruct visual stimuli \citep{scotti2023reconstructing, wang2024decoding, xia2024dream}. Some studies incorporate auxiliary semantic information from language models (LMs)~\citep{lin2022mind, quan2024psychometry}, but still rely on vision-based representations as the core latent space. In contrast, we question whether matching the modality of visual stimuli is truly essential for reconstruction.
%\yy{I think first say what other papers did and then throw out the question sounds more logically.}
%Neuroscience shows that regions near the visual cortex border act as convergence zones where visual input transitions into an amodal semantic system \citep{popham2021visual, tang2023semantic}. However, it remains unclear whether this semantic representation aligns more closely with the latent space of vision models or alternative representations. Although some methods incorporate semantic information from language models \citep{lin2022mind, quan2024psychometry}, such information is often auxiliary, with most approaches relying on direct fMRI-to-vision latent mappings \citep{scotti2023reconstructing, wang2024decoding, xia2024dream}. 
In addition, prior work suffers from limited reconstruction quality due to a unified hidden representation that conflates objects and their attributes, often causing object detection errors, e.g., generating a tiger instead of a gray, tiger-striped cat (Section~\ref{app:common_errors} of the Appendix). This reflects a fundamental mismatch with human visual processing, which is object-centric and compositional rather than holistic \citep{marr1980visual, bracci2023understanding}. Overcoming this limitation calls for generative models that explicitly capture the compositional structure of human perception.
%Yet, it remains an open question which object attributes—such as location, geometry, or semantic category—the brain emphasizes during visual representation.
% \yy{Logic a bit wierd here. seems like you are conflicting with what you said before.} \yy{you should first say the current problem at high-level, e.g. poor reconstruction quality. misattributaion error is just an example of this poor reconstruction quality.}
% The fundamental challenge lies in the synthesis process itself: reconstructed images often suffer from generative models misattributing visual properties—for instance, generating a tiger instead of a gray tiger-striped cat (Appendix~\ref{app:common_errors}). This reflects a mismatch with human visual processing, which is object-centric and compositional rather than holistic \citep{marr1980visual, bracci2023understanding}.\yy{why this reflect you should do object-centric and compositional generation? logic gap.} Addressing this gap requires generative models that better capture the object-based structure of human perception. Yet, it remains an open question which object attributes—such as location, geometry, or semantic category—the brain emphasizes during visual representation.\yy{this question also comes from nowhere. you were mentioning that object detection error. why do you suddenly mention object attributes? }


To address these issues, we propose \modelname, a model that {\textbf{P}}rojects fM{\textbf{R}}I s{\textbf{I}}gnals into a \textbf{S}tructured text space as an inter{\textbf{M}}ediate representation for image reconstruction. To identify the most effective intermediate space, we compare fMRI signals with representations from pre-trained vision, language, and vision–language models using established metrics \citep{wang2020finding, murphy2024correcting, keskar2016large}. Unlike prior work assuming vision-based representations are essential, our first finding (\Cref{sec:findings}) shows that fMRI signals align more closely with the text space of an LM, motivating the use of solely text as a bridge for reconstruction. Building on this, our second finding reveals that reconstruction quality improves when the text and the generative model are adapted to capture the compositional and relational nature of visual stimuli—encompassing objects, their attributes, and their relationships. Guided by these insights, we develop two core modules: an object-centric diffusion module that adapts the diffusion model to generate images by composing individual objects, and an attribute/relationship search module that uses a vision–language model (VLM) to automatically identify object attributes and relationships aligned with neural activity, providing structured guidance for reconstruction.
%\yy{you should highlight the difference with the previous works here. something like, unlike prior work that .../contrast to common beliefs that..., our analysis surprisingly show that ...}
% To further enhance reconstruction quality and mitigate object detection errors, we introduce structured texts that encode objects, attributes, and their relationships, inspired by the compositional nature of human perception. Concretely, our framework follows a two-stage process: (1) we adapt the generation process of pre-trained diffusion models to enable object-centric image synthesis, better reflecting the compositional nature of human visual cognition, and (2) we develop a vision–language model–assisted method to automatically discover which object relations are most semantically meaningful and effectively represented in the brain during visual processing. 
Our contributions are summarized as follows:
% To address the problem of misattribution, we constrain the format of the textual descriptions used to guide generation, ensuring that generative models can capture the object-based structure of human perception. Specifically, we design a structured text format in which each description consists of objects and their attributes, spanning multiple levels of detail—from high-level semantics such as object categories and actions to low-level perceptual attributes such as spatial position and color. Because the most effective attributes should be those that are both semantically meaningful and reliably represented in the brain, we further introduce an automatic attribute discovery method. Starting from a candidate set of attributes, we employ a vision–language model (VLM) to iteratively refine and expand this set, selecting those that exhibit strong alignment with the subject’s fMRI responses and enhance reconstruction quality. Thus, structured object-centric text serves as an intermediate representation, providing a principled bridge between fMRI signals and generative models for image reconstruction. \yy{too detail and lack context. why do people need to know the search detail. how does that help with your claim.}
% To address the challenge of misattributing, we adapt the generation process of pre-trained diffusion models to enable object-centric generation, better reflecting the compositional nature of human visual cognition.
% Finally, to identify which object attributes are most semantically meaningful and effectively represented in the brain during visual processing, we propose an automatic attribute discovery method. 
% Beginning with a set of candidate attributes, we use a VLM to iteratively refine and expand this set, selecting attributes that exhibit strong alignment with the subject’s fMRI responses and enhance reconstruction quality.
% leverage the extensive vocabulary of large language models \citep{yao2024survey, naveed2023comprehensive}. 
% Instead of manually specifying candidate attributes, we utilize prompt optimization \citep{pryzant2023automatic} to explore and identify descriptive words and phrases from this vast vocabulary. This approach allows the model to automatically discover and emphasize the attributes that exhibit the strongest alignment with the subject's fMRI responses, thereby uncovering the most relevant object attributes for visual understanding.
\begin{itemize}[noitemsep,topsep=0pt]
    \item \textbf{Novel Findings:} To our knowledge, we are the first to show that accurate visual stimuli reconstruction can be achieved without image-based latent representations, with LM text space effectively bridging brain activity and generative models. Furthermore, we find that adapting this text space and the generative model to capture the compositional and relational nature of visual images further improves reconstruction quality.
    %In particular, structured textual descriptions derived from LM space serve as effective guidance for decoding visual stimuli.
    
    %We show that fMRI signals align more closely with the semantic embeddings of language models than with the perceptual features of vision models, suggesting that the brain encodes visual information in an abstract, language-like form.
    \item \textbf{Novel Framework:} 
    % Motivated by our empirical findings,
    % %\yy{if you use this word, make sure you mark it in the previous text} 
    % we propose a new fMRI-to-Image reconstruction framework, which enables object-centric reconstruction by adapting diffusion models and automatically discovers brain-aligned object relations and relations using vision–language models.
    Motivated by our empirical findings, we introduce a new fMRI-to-image reconstruction framework that adapts diffusion models for object-centric generation and leverages VLMs to automatically identify brain-aligned object attributes and relationships that can optimally guide the reconstruction.
    \item \textbf{Comprehensive Experiments:} Extensive evaluations on real-world fMRI datasets demonstrate that our method achieves up to an $6\%$ reduction in perceptual loss compared to state-of-the-art models, highlighting the effectiveness of our framework.
\end{itemize}







% to address misattribute, constrain text-> text object -> object span multi details (optimal attributes)-> strctual text
% strctual text -> t5 . combined with learned t5 + object diffusion -> recon